\documentclass[12pt,a4paper]{article}
\usepackage[margin=18mm]{geometry}
\usepackage{fontspec,polyglossia}
\setmainlanguage{ukrainian}
\setmainfont{Times New Roman}
\newfontfamily\cyrillicfont{Times New Roman}
\usepackage{amsmath,array,xcolor}
\setlength{\parindent}{0pt}
\setlength{\parskip}{6pt}
\setlength{\tabcolsep}{6pt}
\renewcommand{\arraystretch}{1.2}
\newenvironment{solutiontable}{\small\begin{tabular}{|>{\raggedright\arraybackslash}p{0.22\linewidth}|>{\raggedright\arraybackslash}p{0.20\linewidth}|>{\raggedright\arraybackslash}p{\dimexpr0.58\linewidth-6\tabcolsep-4\arrayrulewidth\relax}|}\hline
\textbf{Дано}&\textbf{СІ}&\textbf{Розв'язання}\\\hline}{\\\hline\end{tabular}\par}
\newenvironment{solutionwide}{\small\begin{tabular}{|>{\raggedright\arraybackslash}p{0.18\linewidth}|>{\raggedright\arraybackslash}p{0.14\linewidth}|>{\raggedright\arraybackslash}p{\dimexpr0.68\linewidth-6\tabcolsep-4\arrayrulewidth\relax}|}\hline
\textbf{Дано}&\textbf{СІ}&\textbf{Розв'язання}\\\hline}{\\\hline\end{tabular}\par}
\newcommand{\target}[1]{\par\rule{\linewidth}{0.3pt}\par $#1-?$}
\newcommand{\why}[1]{\textbf{\textit{#1}}\par}

\newenvironment{solutiontwo}{\small\begin{tabular}{|>{\raggedright\arraybackslash}p{0.27\linewidth}|>{\raggedright\arraybackslash}p{\dimexpr0.73\linewidth-4\tabcolsep-3\arrayrulewidth\relax}|}\hline
\textbf{Дано}&\textbf{Розв'язання}\\\hline}{\\\hline\end{tabular}\par}
\begin{document}
\begin{center}{\LARGE Механічний рух}\end{center}
\section*{Прямолінійний рух}
\subsection*{Задача 1.28. Обгін поїздів}
\textbf{Умова.} Два поїзди однакової довжини їдуть паралельними коліями в одному напрямку. Швидкість першого поїзда дорівнює 120 км/год, а швидкість другого --- 66 км/год. Перший поїзд обганяє другий протягом 1 хв. Чому дорівнює довжина кожного поїзда?

\textbf{Переведення швидкостей у СІ.} Оскільки $1\,\text{км}=1000\,\text{м}$, а $1\,\text{год}=3600\,\text{с}$, то
\[
1\,\frac{\text{км}}{\text{год}}=\frac{1000\,\text{м}}{3600\,\text{с}}=\frac{1}{3{,}6}\,\frac{\text{м}}{\text{с}}.
\]
Отже, щоб перевести км/год у м/с, числове значення ділимо на $3{,}6$:
\[
v_1=120\cdot\frac{1000}{3600}=\frac{100}{3}\,\text{м/с},\qquad
v_2=66\cdot\frac{1000}{3600}=\frac{55}{3}\,\text{м/с}.
\]
\begin{solutiontable}
$v_1=120\,\text{км/год}$\par
$v_2=66\,\text{км/год}$\par
$t=1\,\text{хв}$\par
$l_1=l_2=l$
\target{l}
&
$v_1=\dfrac{100}{3}\,\text{м/с}$\par
$v_2=\dfrac{55}{3}\,\text{м/с}$\par
$t=60\,\text{с}$
&
\why{Розглянемо рух першого поїзда відносно другого.}
Оскільки поїзди рухаються в одному напрямку, відносна швидкість дорівнює різниці швидкостей:
\[v_{\text{відн}}=v_1-v_2.\]
\why{Для повного обгону треба подолати суму довжин обох поїздів.}
Від початку обгону, коли голова першого поїзда зрівнялася з хвостом другого, до завершення, коли хвіст першого зрівнявся з головою другого, відносний шлях становить
\[s_{\text{відн}}=l_1+l_2=2l.\]
Тому
\[(v_1-v_2)t=2l,\]
\[l=\frac{(v_1-v_2)t}{2}.\]
Обчислимо відносну швидкість:
\[v_{\text{відн}}=\frac{120-66}{3{,}6}=15\,\text{м/с}.\]
\[l=\frac{15\cdot60}{2}=450\,\text{м}.\]
\textbf{Перевірка:} $2\cdot450/15=60\,\text{с}$.
\par\textbf{Відповідь:} довжина кожного поїзда $450\,\text{м}$.
\end{solutiontable}
\newpage
\subsection*{Задача 1.51. Відстань між автомобілями}
\textbf{Умова.} Із двох міст, що розміщені на відстані 300 км один від одного, прямою дорогою одночасно назустріч один одному виїхали вантажівка зі швидкістю 40 км/год і таксі зі швидкістю 60 км/год. Через скільки часу після виїзду відстань між автомобілями дорівнюватиме 100 км?

\begin{solutiontwo}
$s_{\text{міст}} =300\,\text{км}$\par
$s_{\text{зустрічі}}=100\,\text{км}$\par
$v_1=40\,\text{км/год}$\par
$v_2=60\,\text{км/год}$
\target{t_1}\par $t_2-?$

&

Позначимо відстань між містами через $s_{\text{міст}}$, а задану відстань між автомобілями --- через $s_{\text{зустрічі}}$.\par
\why{За зустрічного руху швидкість зближення $v_{\text{збл}}$ дорівнює сумі швидкостей.}
\[v_{\text{збл}} =v_1+v_2.\]
До зустрічі відстань $s_{\text{зустрічі}}=s_{\text{міст}} -v_{\text{збл}} t$, тому
\[t_1=\frac{s_{\text{міст}} -s_{\text{зустрічі}}}{v_1+v_2}.\]
Після зустрічі автомобілі віддаляються, і $s_{\text{зустрічі}}=v_{\text{збл}} t-s_{\text{міст}} $, звідки
\[t_2=\frac{s_{\text{міст}} +s_{\text{зустрічі}}}{v_1+v_2}.\]
Обчислюємо в км та годинах:
\[v_{\text{збл}} =40+60=100\,\text{км/год},\]
\[t_1=\frac{300-100}{100}=2\,\text{год},\]
\[t_2=\frac{300+100}{100}=4\,\text{год}.\]
Зустріч відбудеться через $300/100=3\,\text{год}$. Другий результат передбачає, що автомобілі продовжують рух після зустрічі.
\par\textbf{Відповідь:} через 2 год і через 4 год після виїзду.
\end{solutiontwo}
\newpage
\subsection*{Задача 1.52. Мишко та Сашко}
\textbf{Умова.} Мишко вийшов із селища і пішов прямою дорогою зі швидкістю 5 км/год. Через годину з цього селища в тому самому напрямку виїхав велосипедом Сашко і поїхав зі швидкістю 15 км/год. Через скільки часу після виїзду відстань між ними дорівнює 2,5 км?

\begin{solutiontwo}
$v_{\text{М}}=5\,\text{км/год}$\par
$v_{\text{С}}=15\,\text{км/год}$\par
$t_{\text{зат}}=1\,\text{год}$\par
$s_{\text{зустрічі}}=2{,}5\,\text{км}$
\target{t_1}\par $t_2-?$

&

\why{Відлічуємо час від виїзду Сашка.}
Позначимо початкову відстань до Мишка через $s_{\text{поч}}$, а задану відстань між хлопцями --- через $s_{\text{зустрічі}}$. За годину Мишко відійшов на відстань
\[s_{\text{поч}} =v_{\text{М}}t_{\text{зат}}.\]
Сашко скорочує відстань зі швидкістю зближення $v_{\text{збл}}$:
\[v_{\text{збл}} =v_{\text{С}}-v_{\text{М}}.\]
До наздоганяння $s_{\text{зустрічі}}=s_{\text{поч}} -v_{\text{збл}} t_1$, а після нього $s_{\text{зустрічі}}=v_{\text{збл}} t_2-s_{\text{поч}} $. Тому
\[t_1=\frac{s_{\text{поч}} -s_{\text{зустрічі}}}{v_{\text{збл}} },\qquad t_2=\frac{s_{\text{поч}} +s_{\text{зустрічі}}}{v_{\text{збл}} }.\]
Підставляємо числа в км та годинах:
\[s_{\text{поч}} =5\cdot1=5\,\text{км},\quad v_{\text{збл}} =15-5=10\,\text{км/год}.\]
\[t_1=\frac{5-2{,}5}{10}=0{,}25\,\text{год},\]
\[t_2=\frac{5+2{,}5}{10}=0{,}75\,\text{год}.\]
Переводимо в хвилини: $0{,}25\cdot60=15$, $0{,}75\cdot60=45$.
\par\textbf{Відповідь:} через 15 хв і через 45 хв після виїзду Сашка.
\end{solutiontwo}
\newpage
\subsection*{Задача 1.64. Середня швидкість Петрика}
\textbf{Умова.} Петрик пройшов 10 км за 2 год, а потім проїхав велосипедом 20 км за 1 год. Чому дорівнює середня швидкість Петрика?

\begin{solutiontwo}
$s_1=10\,\text{км}$\par
$t_1=2\,\text{год}$\par
$s_2=20\,\text{км}$\par
$t_2=1\,\text{год}$
\target{v_{\text{сер}}}

&

\why{Середня швидкість дорівнює всьому шляху, поділеному на весь час руху.}
\[v_{\text{сер}}=\frac{s}{t}.\]
Для двох ділянок $s=s_1+s_2$, $t=t_1+t_2$, отже
\[v_{\text{сер}}=\frac{s_1+s_2}{t_1+t_2}.\]
Обчислимо в кілометрах і годинах:
\[v_{\text{сер}}=\frac{10+20}{2+1}=10\,\text{км/год}.\]
Середнє арифметичне швидкостей ділянок тут не підходить, бо час руху на них різний.
\par\textbf{Відповідь:} $10\,\text{км/год}$.
\end{solutiontwo}
\newpage
\subsection*{Задача О-14. Середня швидкість автомобіля}
\textbf{Умова.} Автомобіль проїхав половину шляху зі швидкістю 60 км/год; половину решти часу він їхав зі швидкістю 15 км/год, а останню ділянку --- зі швидкістю 45 км/год. Яка середня швидкість на всьому шляху?

\begin{solutiontwo}
$v_1=60\,\text{км/год}$\par
$v_2=15\,\text{км/год}$\par
$v_3=45\,\text{км/год}$\par
$s_1=s/2$\par
$t_2=t_3$
\target{v_{\text{сер}}}

&

\setlength{\abovedisplayskip}{5pt}\setlength{\belowdisplayskip}{5pt}
\why{Крок 1. Позначимо весь шлях через $s$.}
Перша ділянка становить половину шляху: $s_1=s/2$. Із формули $v_1=s_1/t_1$ виражаємо час: множимо на $t_1$ і ділимо на $v_1$.
\[v_1t_1=s_1,\qquad t_1=\frac{s_1}{v_1}=\frac{s}{2v_1}.\]
\why{Крок 2. Запишемо шляхи на двох останніх ділянках.}
Після першої ділянки залишилося $s-s/2=s/2$. За умовою половину решти часу автомобіль їде з $v_2$, іншу половину --- з $v_3$. Тому $t_3=t_2$.
\[s_2=v_2t_2,\qquad s_3=v_3t_3=v_3t_2.\]
Сума цих шляхів дорівнює $s/2$:
\[\frac{s}{2}=v_2t_2+v_3t_2=(v_2+v_3)t_2.\]
\why{Крок 3. Виразимо час другої половини шляху.}
Ділимо обидві частини рівняння на $v_2+v_3$:
\[t_2=\frac{s}{2(v_2+v_3)}.\]
Оскільки таких однакових проміжків часу два,
\[t_{23}=t_2+t_3=2t_2=\frac{2s}{2(v_2+v_3)}=\frac{s}{v_2+v_3}.\]
\why{Крок 4. Знайдемо весь час руху.}
\[t=t_1+t_{23}=\frac{s}{2v_1}+\frac{s}{v_2+v_3}.\]
Спільний знаменник --- $2v_1(v_2+v_3)$. Перший чисельник множимо на $v_2+v_3$, другий --- на $2v_1$:
\[t=\frac{s(v_2+v_3)+2v_1s}{2v_1(v_2+v_3)}=\frac{s(v_2+v_3+2v_1)}{2v_1(v_2+v_3)}.\]
\why{Крок 5. Поділимо весь шлях на весь час.}
Ділення на дріб замінюємо множенням на обернений дріб. Додатний шлях $s$ скорочується:
\[
\begin{aligned}
v_{\text{сер}}&=\frac{s}{t}=s\cdot\frac{2v_1(v_2+v_3)}{s(v_2+v_3+2v_1)}\\
&=\frac{2v_1(v_2+v_3)}{v_2+v_3+2v_1}.
\end{aligned}
\]
\why{Крок 6. Підставимо швидкості в км/год.}
\[v_{\text{сер}}=\frac{2\cdot60\cdot(15+45)}{15+45+2\cdot60}=\frac{7200}{180}=40\,\text{км/год}.\]
\textbf{Відповідь:} $40\,\text{км/год}$.
\end{solutiontwo}
\newpage
\subsection*{Задача О-15. Велосипедист і пішохід}
\textbf{Умова.} Велосипедист половину часу всього руху їхав зі швидкістю 20 км/год, половину решти шляху зі швидкістю 12 км/год, а останню ділянку --- йшов зі швидкістю 6 км/год. Яка середня швидкість на всьому шляху?

\begin{solutiontwo}
$v_1=20\,\text{км/год}$\par
$v_2=12\,\text{км/год}$\par
$v_3=6\,\text{км/год}$\par
$t_1=t /2$\par
$s_2=s_3$
\target{v_{\text{сер}}}

&

\why{На двох останніх ділянках шляхи однакові, а час руху різний.}
Позначимо весь шлях через $s$, а довжину другої ділянки --- через $s_2$. Третя ділянка має таку саму довжину: $s_3=s_2$. Тоді
\[t_2=\frac{s_2}{v_2},\qquad t_3=\frac{s_2}{v_3}.\]
Середня швидкість на решті шляху:
\[
\begin{aligned}
v_{23}&=\frac{2s_2}{s_2/v_2+s_2/v_3}\\
&=\frac{2v_2v_3}{v_2+v_3}.
\end{aligned}
\]
\why{Перша частина руху та решта руху займають однаковий час $t /2$.}
Загальний шлях і середня швидкість:
\[s=v_1\frac{t }{2}+v_{23}\frac{t }{2},\]
\[v_{\text{сер}}=\frac{s}{t }=\frac{v_1+v_{23}}{2}.\]
Підставляємо числа:
\[v_{23}=\frac{2\cdot12\cdot6}{12+6}=8\,\text{км/год},\]
\[v_{\text{сер}}=\frac{20+8}{2}=14\,\text{км/год}.\]
\textbf{Відповідь:} $14\,\text{км/год}$.
\end{solutiontwo}
\newpage
\subsection*{Задача О-16. Мандрівник на коні та віслюку}
\textbf{Умова.} Мандрівник їхав спочатку на коні, а потім на віслюку. Яку частину шляху і яку частину всього часу руху він їхав на коні, якщо середня швидкість мандрівника, як виявилося, дорівнювала 12 км/год, швидкість їзди на коні 30 км/год, а на віслюку --- 6 км/год?

\begin{solutiontwo}
$v_{\text{сер}}=12\,\text{км/год}$\par
$v_{\text{к}}=30\,\text{км/год}$\par
$v_{\text{в}}=6\,\text{км/год}$
\target{t_{\text{к}}/t }\par
$ s_{\text{к}}/s-?$

&

\why{Виразимо весь шлях через час руху на коні та на віслюку.}
Нехай $t $ --- увесь час, $t_{\text{к}}$ --- час на коні. Тоді на віслюку мандрівник їхав $t -t_{\text{к}}$:
\[s=v_{\text{к}}t_{\text{к}}+v_{\text{в}}(t -t_{\text{к}}).\]
Водночас $s=v_{\text{сер}}t $. Отже,
\[(v_{\text{сер}}-v_{\text{в}})t =(v_{\text{к}}-v_{\text{в}})t_{\text{к}}.\]
Частина часу, проведена на коні:
\[\frac{t_{\text{к}}}{t }=\frac{v_{\text{сер}}-v_{\text{в}}}{v_{\text{к}}-v_{\text{в}}}.\]
\why{Частку шляху знаходимо як відношення шляху на коні до всього шляху.}
\[\frac{s_{\text{к}}}{s}=\frac{v_{\text{к}}t_{\text{к}}}{v_{\text{сер}}t }=\frac{v_{\text{к}}}{v_{\text{сер}}}\cdot\frac{t_{\text{к}}}{t }.\]
Підставляємо числа:
\[\frac{t_{\text{к}}}{t }=\frac{12-6}{30-6}=\frac{6}{24}=\frac14,\]
\[\frac{s_{\text{к}}}{s}=\frac{30}{12}\cdot\frac14=\frac58.\]
\textbf{Відповідь:} на коні мандрівник проїхав $\frac58$ усього шляху та витратив $\frac14$ всього часу руху.
\end{solutiontwo}
\newpage
\subsection*{Задача О-17. П'ять кілометрів спортсмена}
\textbf{Умова.} Спортсмен подолав дистанцію 5 км. Перший кілометр він пробіг за 3 хв, а на кожний наступний кілометр у нього йшло на $t$ секунд більше, ніж на попередній. Знайдіть $t$, якщо відомо, що середня швидкість на всьому шляху виявилася такою, як коли б спортсмен пробігав кожний кілометр за 3 хв 12 с.

\begin{solutionwide}
$s=5\,\text{км}$\par
$s_1=1\,\text{км}$\par
$t_1=3\,\text{хв}$\par
$t_0=3\,\text{хв}\,12\,\text{с}$
\target{\Delta t}

&

$t_1=3\cdot60$\par
$\phantom{t_1}=180\,\text{с}$\par
$t_0=3\cdot60+12$\par
$\phantom{t_0}=192\,\text{с}$

&

\setlength{\abovedisplayskip}{4pt}\setlength{\belowdisplayskip}{4pt}
\why{Крок 1. Визначимо кількість однакових ділянок.}
Дистанція становить 5 км, довжина кожної ділянки --- 1 км. Отже, маємо п’ять однакових за довжиною ділянок: $5:1=5$.
Позначення: $t_1$ --- час першого кілометра; $t_0$ --- час одного кілометра за рівномірного бігу; $\Delta t$ --- шукане збільшення часу (у самій умові воно позначене буквою $t$); увесь час руху позначимо $t$.
\par
\why{Крок 2. Знайдемо весь час за заданою середньою швидкістю.}
Якби кожний кілометр долали за $t_0$, швидкість дорівнювала б $v_0=s_1/t_0$. За умовою фактична середня швидкість така сама:
\[\frac{s}{t }=\frac{s_1}{t_0}.\]
Множимо обидві частини на $t t_0$:
\[st_0=s_1 t .\]
Ділимо на $s_1$ і враховуємо $s/s_1=5$:
\[t =\frac{s}{s_1}t_0=5t_0.\]
\why{Крок 3. Запишемо час кожного кілометра.}
Перший: $t_1$. Другий: $t_1+\Delta t$.
Третій: $(t_1+\Delta t)+\Delta t=t_1+2\Delta t$.
Четвертий: $(t_1+2\Delta t)+\Delta t=t_1+3\Delta t$.
П'ятий: $(t_1+3\Delta t)+\Delta t=t_1+4\Delta t$.
\par
\why{Крок 4. Додамо п'ять проміжків часу.}
\[
\begin{aligned}
t ={}&t_1+(t_1+\Delta t)+(t_1+2\Delta t)\\
&+(t_1+3\Delta t)+(t_1+4\Delta t).
\end{aligned}
\]
Час $t_1$ повторюється п'ять разів, а додатковий час становить $(1+2+3+4)\Delta t=10\Delta t$. Тому
\[t =5t_1+10\Delta t.\]
\why{Крок 5. Прирівняємо вирази для всього часу.}
\[5t_0=5t_1+10\Delta t.\]
Віднімаємо $5t_1$ від обох частин:
\[5t_0-5t_1=10\Delta t.\]
Виносимо 5 за дужки й ділимо на 10:
\[\Delta t=\frac{5(t_0-t_1)}{10}=\frac{t_0-t_1}{2}.\]
\why{Крок 6. Підставимо час у секундах.}
\[\Delta t=\frac{192-180}{2}=\frac{12}{2}=6\,\text{с}.\]
\textbf{Перевірка:} $180+186+192+198+204=960\,\text{с}$; також $5\cdot192=960\,\text{с}$.
\par\textbf{Відповідь:} $\Delta t=6\,\text{с}$.
\end{solutionwide}
\newpage
\subsection*{Задача 2. Повернення додому}
\textbf{Умова.} Хлопчик іде до школи зі швидкістю 3 км/год і розраховує встигнути до початку уроку. Пройшовши половину шляху, він згадує, що забув зачинити вхідні двері. З якою швидкістю повинен бігти хлопчик, щоб повернутись додому, зачинити двері і не запізнитись до школи? (Час, витрачений на зачинення дверей, не враховувати.)

\begin{solutiontwo}
$v_0=3\,\text{км/год}$\par
$s_1=s/2$\par
$t_{\text{двері}}=0$
\target{v}

&

\why{Порівняємо запланований решту шляху з тим, який тепер треба пробігти.}
Нехай $s$ --- відстань від дому до школи. Вважаємо, що за початковим планом хлопчик мав прийти саме до початку уроку.
Після проходження половини шляху залишився час
\[t_{\text{зал}}=\frac{s/2}{v_0}=\frac{s}{2v_0}.\]
\why{Тепер хлопчик має повернутися пів шляху додому й пробігти весь шлях до школи.}
\[s_{\text{біг}}=\frac{s}{2}+s=\frac{3s}{2}.\]
Найменша стала швидкість бігу:
\[
\begin{aligned}
v&=\frac{s_{\text{біг}}}{t_{\text{зал}}}\\
&=\frac{3s/2}{s/(2v_0)}=3v_0.
\end{aligned}
\]
Підставляємо число:
\[v=3\cdot3=9\,\text{км/год}.\]
\textbf{Відповідь:} не менше $9\,\text{км/год}$ за відсутності початкового запасу часу.
\end{solutiontwo}
\newpage
\subsection*{Задача 1. Туристи: пішки, автобусом та електричкою}
\textbf{Умова (на фото позначено «9 клас»).} Шлях від місця відпочинку до міста туристи подолали із середньою швидкістю 32 км/год частково пішки, частково автобусом, а потім електричкою. Визначити, з якою швидкістю пройдено кожен із відрізків шляху, якщо їхні довжини відносяться як $1:4:45$, а відповідні інтервали часу --- як $4:1:20$.

\begin{solutiontwo}
$v_{\text{сер}}=32\,\text{км/год}$\par
$s_1:s_2:s_3$\par
$\phantom{s_1}=1:4:45$\par
$t_1:t_2:t_3$\par
$\phantom{t_1}=4:1:20$
\target{v_1}\par
$v_2-?$\par
$v_3-?$

&

\why{Запишемо довжини та часи через одну спільну частину кожного відношення.}
Нехай $a$ --- одна частина шляху, $b$ --- одна частина часу. Тоді
\[s_1=a,\quad s_2=4a,\quad s_3=45a;\]
\[t_1=4b,\quad t_2=b,\quad t_3=20b.\]
Увесь шлях і весь час:
\[s=(1+4+45)a=50a,\]
\[t =(4+1+20)b=25b.\]
\why{Середня швидкість дає змогу знайти відношення $a/b$.}
\[v_{\text{сер}}=\frac{s}{t }=\frac{50a}{25b}=\frac{2a}{b},\]
\[\frac{a}{b}=\frac{v_{\text{сер}}}{2}=\frac{32}{2}=16\,\text{км/год}.\]
\why{Для кожної ділянки ділимо її шлях на відповідний час.}
\[v_1=\frac{a}{4b}=\frac14\cdot16=4\,\text{км/год},\]
\[v_2=\frac{4a}{b}=4\cdot16=64\,\text{км/год},\]
\[v_3=\frac{45a}{20b}=\frac{45}{20}\cdot16=36\,\text{км/год}.\]
\textbf{Відповідь:} пішки --- $4\,\text{км/год}$; автобусом --- $64\,\text{км/год}$; електричкою --- $36\,\text{км/год}$.
\end{solutiontwo}
\newpage
\subsection*{Задача 8. Середня швидкість поїзда}
\textbf{Умова.} Визначити середню швидкість поїзда, якщо на проходження окремих ділянок, довжини яких відносяться як $1:3:4:2$, потрібно було затратити часу у відношенні $2:4:3:1$ і на останній ділянці швидкість дорівнювала 80 км/год.

\begin{solutiontwo}
$s_1:s_2:s_3:s_4$\par
$\phantom{s_1}=1:3:4:2$\par
$t_1:t_2:t_3:t_4$\par
$\phantom{t_1}=2:4:3:1$\par
$v_4=80\,\text{км/год}$
\target{v_{\text{сер}}}
&
\why{Крок 1. Виразимо всі ділянки через довжину першої ділянки $s$.}
З указаного відношення довжин:
\[s_1=s,\quad s_2=3s,\quad s_3=4s,\quad s_4=2s.\]
Тому весь шлях дорівнює
\[s_{\text{заг}}=s+3s+4s+2s=10s.\]
\why{Крок 2. Виразимо часи через час останньої ділянки $t_4$.}
Останній проміжок часу становить одну частину відношення. Отже,
\[t_1=2t_4,\quad t_2=4t_4,\quad t_3=3t_4.\]
Увесь час руху:
\[t=2t_4+4t_4+3t_4+t_4=10t_4.\]
\why{Крок 3. Скористаємося швидкістю на останній ділянці.}
\[v_4=\frac{s_4}{t_4}=\frac{2s}{t_4}.\]
Ділимо обидві частини на 2:
\[\frac{s}{t_4}=\frac{v_4}{2}.\]
\why{Крок 4. Ділимо весь шлях на весь час.}
\[v_{\text{сер}}=\frac{s_{\text{заг}}}{t}=\frac{10s}{10t_4}=\frac{s}{t_4}=\frac{v_4}{2}.\]
Підставляємо число:
\[v_{\text{сер}}=\frac{80}{2}=40\,\text{км/год}.\]
\textbf{Відповідь:} $40\,\text{км/год}$.
\end{solutiontwo}
\newpage
\subsection*{Задача 1.58. Рух ескалатором}
\textbf{Умова.} Коли Петрик стоїть на сходинці ескалатора, він спускається за 1 хв. Коли Петрик іде по ескалатору, що рухається вниз, він спускається за 30 с. Скільки часу Петрик спускатиметься, якщо йтиме удвічі швидше?

\begin{solutiontable}
$t_0=1\,\text{хв}$\par
$t_1=30\,\text{с}$\par
$v'_\text{п}=2v_\text{п}$
\target{t_2}
&
$t_0=1\cdot60=60\,\text{с}$
&
\why{Швидкості ескалатора і Петрика відносно сходинок додаються.}
Нехай $l$ --- довжина спуску, $v_{\text{е}}$ --- швидкість ескалатора, $v_{\text{п}}$ --- швидкість ходьби Петрика відносно ескалатора.
\[v_{\text{е}}=\frac{l}{t_0},\qquad v_{\text{е}}+v_{\text{п}}=\frac{l}{t_1}.\]
Віднімаємо перше рівняння від другого:
\[v_{\text{п}}=\frac{l}{t_1}-\frac{l}{t_0}.\]
\why{У третьому випадку подвоюється швидкість ходьби, а швидкість ескалатора стала.}
\[
\begin{aligned}
v_2&=v_{\text{е}}+2v_{\text{п}}\\
&=\frac{l}{t_0}+2\left(\frac{l}{t_1}-\frac{l}{t_0}\right)\\
&=l\left(\frac{2}{t_1}-\frac{1}{t_0}\right).
\end{aligned}
\]
Тоді
\[t_2=\frac{l}{v_2}=\frac{t_0t_1}{2t_0-t_1}.\]
\[t_2=\frac{60\cdot30}{2\cdot60-30}=20\,\text{с}.\]
\textbf{Відповідь:} $20\,\text{с}$.
\end{solutiontable}
\newpage
\section*{Рух за течією та проти течії}
\begingroup\color{red!80!black}
\subsection*{Правило для руху за течією та проти течії}
$v_{\text{ч}}$ --- власна швидкість човна відносно води; $v_{\text{р}}$ --- швидкість течії річки; $v_{\text{за}}$ і $v_{\text{проти}}$ --- швидкості човна відносно берега. Для руху проти течії потрібно $v_{\text{ч}}>v_{\text{р}}$.
\[
\begin{cases}
v_{\text{ч}}+v_{\text{р}}=v_{\text{за}},\\
v_{\text{ч}}-v_{\text{р}}=v_{\text{проти}}.
\end{cases}
\]
\textbf{Додаємо рівняння почленно:} додаємо ліві частини та окремо праві. Доданки $+v_{\text{р}}$ і $-v_{\text{р}}$ взаємно знищуються:
\[
(v_{\text{ч}}+v_{\text{р}})+(v_{\text{ч}}-v_{\text{р}})=v_{\text{за}}+v_{\text{проти}},
\]
\[2v_{\text{ч}}=v_{\text{за}}+v_{\text{проти}},\qquad
v_{\text{ч}}=\frac{v_{\text{за}}+v_{\text{проти}}}{2}.
\]
\textbf{Віднімаємо друге рівняння від першого:} знак мінус перед дужками змінює знаки обох доданків у дужках:
\[
(v_{\text{ч}}+v_{\text{р}})-(v_{\text{ч}}-v_{\text{р}})=v_{\text{за}}-v_{\text{проти}},
\]
\[2v_{\text{р}}=v_{\text{за}}-v_{\text{проти}},\qquad
v_{\text{р}}=\frac{v_{\text{за}}-v_{\text{проти}}}{2}.
\]
\textbf{Якщо відомі власна швидкість і швидкість течії}, обчислюємо окремо:
\[v_{\text{за}}=v_{\text{ч}}+v_{\text{р}},\qquad v_{\text{проти}}=v_{\text{ч}}-v_{\text{р}}.\]
\textbf{Як користуватися правилом.} Додавання двох рівнянь дає власну швидкість човна, віднімання --- швидкість течії. Швидкість за течією знаходимо з першого рівняння, проти течії --- з другого. Це правило можна використовувати в усіх наступних задачах про рух річкою.
\endgroup
\newpage
\subsection*{Рух човна за течією та проти течії річки}
\textbf{Умова (8 клас).} Моторний човен проплив 24 км за течією річки за 1,5 год, а таку саму відстань проти течії --- за 2 год.

Визначте: 1) швидкість човна за течією річки; 2) швидкість човна проти течії річки; 3) власну швидкість човна (у стоячій воді); 4) швидкість течії річки; 5) на скільки хвилин швидше човен подолає відстань 36 км за течією, ніж проти течії.

\begin{solutiontwo}
$s=24\,\text{км}$\par
$t_{\text{з}}=1{,}5\,\text{год}$\par
$t_{\text{п}}=2\,\text{год}$\par
$s_2 =36\,\text{км}$
\target{v_{\text{за}}}\par
$v_{\text{проти}}-?$\par
$v_{\text{ч}}-?$\par
$v_{\text{р}} -?$\par
$\Delta t-?$

&

\why{1--2. Швидкість дорівнює шляху, поділеному на час.}
Індекси «за» і «проти» позначають рух за течією та проти течії:
\[v_{\text{за}}=\frac{s}{t_{\text{з}}},\qquad v_{\text{проти}}=\frac{s}{t_{\text{п}}}.\]
\[v_{\text{за}}=\frac{24}{1{,}5}=16\,\text{км/год},\]
\[v_{\text{проти}}=\frac{24}{2}=12\,\text{км/год}.\]
\why{3--4. За течією швидкості додаються, проти течії --- віднімаються.}
Нехай $v_{\text{ч}}$ --- власна швидкість човна, $v_{\text{р}} $ --- швидкість течії. Тоді
\[
\begin{cases}
v_{\text{ч}}+v_{\text{р}}=v_{\text{за}},\\
v_{\text{ч}}-v_{\text{р}}=v_{\text{проти}}.
\end{cases}
\]

Додавши рівняння, маємо $2v_{\text{ч}}=v_{\text{за}}+v_{\text{проти}}$. Віднявши друге від першого, маємо $2v_{\text{р}} =v_{\text{за}}-v_{\text{проти}}$. Отже,
\[v_{\text{ч}}=\frac{v_{\text{за}}+v_{\text{проти}}}{2}=\frac{16+12}{2}=14\,\text{км/год},\]
\[v_{\text{р}} =\frac{v_{\text{за}}-v_{\text{проти}}}{2}=\frac{16-12}{2}=2\,\text{км/год}.\]
\why{5. Знайдемо різницю часу проходження 36 км.}
\[\Delta t=\frac{s_2 }{v_{\text{проти}}}-\frac{s_2 }{v_{\text{за}}}.\]
\[\Delta t=\frac{36}{12}-\frac{36}{16}=3-2{,}25=0{,}75\,\text{год}.\]
\[\Delta t=0{,}75\cdot60=45\,\text{хв}.\]
\textbf{Відповідь:} 1) $16\,\text{км/год}$; 2) $12\,\text{км/год}$; 3) $14\,\text{км/год}$; 4) $2\,\text{км/год}$; 5) на $45\,\text{хв}$ швидше.
\end{solutiontwo}
\newpage
\subsection*{Додаткова задача. Відстань між пристанями}
\textbf{Умова (підвищена складність).} Човен вирушив із пристані А до пристані Б, розташованої нижче за течією, і одразу повернувся назад. Уся подорож тривала 3,5 год. Визначте відстань між пристанями, якщо швидкості човна залишалися незмінними.

Використайте дані основної задачі: власна швидкість човна $14\,\text{км/год}$, швидкість течії $2\,\text{км/год}$.

\begin{solutiontwo}
$t =3{,}5\,\text{год}$\par
$v_{\text{ч}}=14\,\text{км/год}$\par
$v_{\text{р}}=2\,\text{км/год}$
\target{l}

&

\setlength{\abovedisplayskip}{5pt}\setlength{\belowdisplayskip}{5pt}
\why{Запишемо систему з правила й додамо рівняння.}
\[
\begin{cases}
v_{\text{ч}}+v_{\text{р}}=v_{\text{за}},\\
v_{\text{ч}}-v_{\text{р}}=v_{\text{проти}}.
\end{cases}
\]
\[2v_{\text{ч}}=v_{\text{за}}+v_{\text{проти}}.\]
Отже, сума швидкостей відносно берега дорівнює $2\cdot14=28\,\text{км/год}$.
\why{Знайдемо кожну швидкість окремо з відповідного рівняння системи.}
\[v_{\text{за}}=14+2=16\,\text{км/год},\]
\[v_{\text{проти}}=14-2=12\,\text{км/год}.\]
\why{На обох ділянках відстань однакова, тому весь час дорівнює сумі двох відношень шляху до швидкості.}
Нехай $l$ --- відстань між пристанями. Час руху до Б та назад:
\[t_{\text{з}}=\frac{l}{v_{\text{за}}},\qquad t_{\text{п}}=\frac{l}{v_{\text{проти}}}.\]
Без зупинки в Б:
\[t =t_{\text{з}}+t_{\text{п}}=\frac{l}{v_{\text{за}}}+\frac{l}{v_{\text{проти}}}.\]
\why{Зведемо дроби до спільного знаменника й виразимо відстань.}
\[t =\frac{l(v_{\text{проти}}+v_{\text{за}})}{v_{\text{за}}v_{\text{проти}}},\]
\[l=\frac{t v_{\text{за}}v_{\text{проти}}}{v_{\text{за}}+v_{\text{проти}}}.\]
Підставляємо швидкості в км/год і час у годинах:
\[l=\frac{3{,}5\cdot16\cdot12}{16+12}=\frac{672}{28}=24\,\text{км}.\]
\textbf{Перевірка:}
\[t_{\text{з}}=\frac{24}{16}=1{,}5\,\text{год},\quad t_{\text{п}}=\frac{24}{12}=2\,\text{год},\]
\[t_{\text{з}}+t_{\text{п}}=1{,}5+2=3{,}5\,\text{год}.\]
\textbf{Відповідь:} відстань між пристанями $24\,\text{км}$.
\end{solutiontwo}
\newpage
\subsection*{Задача 1.56. Човен у річці й озері}
\textbf{Умова.} Моторний човен проплив річкою від пристані А до пристані Б і назад за 8 год. Таку саму відстань цей човен проплив в озері за 6 год. Швидкість човна відносно води не змінювалася. У скільки разів швидкість човна відносно води більша за швидкість течії річки?

\begin{solutiontwo}
$t_{\text{р}}=8\,\text{год}$\par
$t_{\text{оз}}=6\,\text{год}$\par
$v_{\text{ч}}$ --- швидкість човна відносно води\par
$v_{\text{р}} $ --- швидкість течії
\target{v_{\text{ч}}/v_{\text{р}}}

&

\setlength{\abovedisplayskip}{5pt}\setlength{\belowdisplayskip}{5pt}
\why{Використаємо систему з правила.}
\[
\begin{cases}
v_{\text{ч}}+v_{\text{р}}=v_{\text{за}},\\
v_{\text{ч}}-v_{\text{р}}=v_{\text{проти}}.
\end{cases}
\]
Додавання дає $v_{\text{за}}+v_{\text{проти}}=2v_{\text{ч}}$. Окремо з рівнянь: $v_{\text{за}}=v_{\text{ч}}+v_{\text{р}}$, $v_{\text{проти}}=v_{\text{ч}}-v_{\text{р}}$.
\why{Запишемо час руху за течією та проти неї.}
Нехай $l$ --- відстань між пристанями. Швидкості відносно берега: $v_{\text{ч}}+v_{\text{р}} $ і $v_{\text{ч}}-v_{\text{р}} $, причому $v_{\text{ч}}>v_{\text{р}} $.
\[
\begin{aligned}
t_{\text{р}}&=\frac{l}{v_{\text{ч}}+v_{\text{р}} }+\frac{l}{v_{\text{ч}}-v_{\text{р}} }\\
&=\frac{l(v_{\text{ч}}-v_{\text{р}} )+l(v_{\text{ч}}+v_{\text{р}} )}{v_{\text{ч}}^2-v_{\text{р}} ^2}\\
&=\frac{2l v_{\text{ч}}}{v_{\text{ч}}^2-v_{\text{р}} ^2}.
\end{aligned}
\]
\why{В озері човен долає такий самий повний шлях $2l$.}
\[t_{\text{оз}}=\frac{2l}{v_{\text{ч}}}.\]
Ділимо час в озері на час у річці:
\[\frac{t_{\text{оз}}}{t_{\text{р}}}=\frac{v_{\text{ч}}^2-v_{\text{р}} ^2}{v_{\text{ч}}^2}=1-\left(\frac{v_{\text{р}} }{v_{\text{ч}}}\right)^2.\]
Звідси
\[\left(\frac{v_{\text{р}} }{v_{\text{ч}}}\right)^2=1-\frac{t_{\text{оз}}}{t_{\text{р}}},\]
Підставляємо час в однакових одиницях:
\[\left(\frac{v_{\text{р}} }{v_{\text{ч}}}\right)^2=1-\frac{6}{8}=\frac14.\]
Отже, $v_{\text{ч}}^2=4v_{\text{р}} ^2$. Оскільки
\[(2v_{\text{р}} )^2=4v_{\text{р}} ^2,\]
а обидві швидкості додатні, маємо $v_{\text{ч}}=2v_{\text{р}} $.
Таким чином,
\[\frac{v_{\text{ч}}}{v_{\text{р}} }=2.\]
\textbf{Відповідь:} у 2 рази.
\end{solutiontwo}
\newpage
\subsection*{Задача 1.57. Час руху човна в озері}
\textbf{Умова.} Човен проплив з А в Б за течією за 2 год, а назад за 8 год. За скільки часу човен пропливе таку саму відстань в озері?

\begin{solutiontwo}
$t_1=2\,\text{год}$\par
$t_2=8\,\text{год}$
\target{t_{\text{оз}}}

&

\why{Власна швидкість човна дорівнює півсумі швидкостей за течією та проти неї.}
Нехай $l=AB$, $v_{\text{ч}}$ --- швидкість човна відносно води, $v_{\text{р}} $ --- швидкість течії.
\[
\begin{cases}
v_{\text{ч}}+v_{\text{р}}=v_{\text{за}},\\
v_{\text{ч}}-v_{\text{р}}=v_{\text{проти}}.
\end{cases}
\]
Окремо швидкості відносно берега:
\[v_{\text{за}}=\frac{l}{t_1},\qquad v_{\text{проти}}=\frac{l}{t_2}.\]

Додаємо рівняння системи: $2v_{\text{ч}}=v_{\text{за}}+v_{\text{проти}}$. Підставляємо знайдені вирази:
\[2v_{\text{ч}}=\frac{l}{t_1}+\frac{l}{t_2},\]
\[v_{\text{ч}}=\frac{l(t_1+t_2)}{2t_1t_2}.\]
\why{Для проходження відстані $AB$ в озері потрібен час $l/v_{\text{ч}}$.}
\[t_{\text{оз}}=\frac{l}{v_{\text{ч}}}=\frac{2t_1t_2}{t_1+t_2}.\]
\[t_{\text{оз}}=\frac{2\cdot2\cdot8}{2+8}=3{,}2\,\text{год}.\]
Дробова частина години:
\[0{,}2\cdot60=12\,\text{хв}.\]
\textbf{Відповідь:} $3\,\text{год}\,12\,\text{хв}$ на відстань $AB$.
\par Якщо розглядати весь шлях туди й назад ($2l$), час в озері дорівнює $6{,}4\,\text{год}=6\,\text{год}\,24\,\text{хв}$.
\end{solutiontwo}
\end{document}
